\section{Universal secant bundles and syzygies}

% archon:covers MR4213770UniversalSecantBundlesAndSyzygiesOfCanonicalCurves.lean MR4213770UniversalSecantBundlesAndSyzygiesOfCanonicalCurves/Basic.lean

\subsection{Introduction}

\begin{theorem}[Generic Green's conjecture for general canonical curves]
  \label{thm:kemeny_generic_green}
  \lean{MR4213770.GenericGreen}
  \uses{lem:kemeny_even_k3_section_to_general_curve, thm:kemeny_odd_voisin}
  % SOURCE: Kemeny, Introduction, Theorem generic-green (read from references/kemeny-universal-secant-bundles-syzygies-canonical-curves-source/kemeny-universal-secant-bundles-syzygies-canonical-curves.tex)
  % SOURCE QUOTE: \begin{thm} \label{generic-green} Let $C$ be a general curve over $\C$ of genus $g=2k$ or $g=2k+1$. Then $\mathrm{K}_{k,1}(C,\omega_C)=0$. \end{thm}
  \textit{Source: Kemeny, Introduction, Theorem generic-green.}
  If \(C\) is a general complex curve of genus \(2k\) or \(2k+1\), then
  \(\K_{k,1}(C,\omega_C)=0\).
\end{theorem}
\begin{proof}
  \uses{lem:kemeny_even_k3_section_to_general_curve, thm:kemeny_odd_voisin}
  For even genus, apply the K3-section transfer from the Picard-rank-one K3
  vanishing proved by the universal secant bundle argument to a general
  canonical curve of genus \(2k\).
  For odd genus, use the Picard-rank-two nodal K3 reduction: the projection
  criterion reduces the desired vanishing to injectivity of the projection map,
  and the odd-genus secant-bundle calculation supplies that injectivity.
\end{proof}

\begin{theorem}[Even-genus geometric syzygy theorem]
  \label{thm:kemeny_even_geometric_syzygy_structure}
  \lean{MR4213770.EvenGeometricSyzygyStructure}
  \uses{cor:kemeny_extremal_syzygy_iso, thm:kemeny_veronese_identification}
  % SOURCE: Kemeny, Introduction, Theorem geo-thm-intro (read from references/kemeny-universal-secant-bundles-syzygies-canonical-curves-source/kemeny-universal-secant-bundles-syzygies-canonical-curves.tex)
  % SOURCE QUOTE: For any nonzero $s \in \mathrm{H}^0(E)$, the space $\mathrm{K}_{k-1,1}(X,L,\mathrm{H}^0(L \otimes I_{Z(s)}))$ is a one-dimensional subspace of $\mathrm{K}_{k-1,1}(X,L)$.
  \textit{Source: Kemeny, Introduction, Theorem geo-thm-intro.}
  Let \(X\) be a complex K3 surface with \(\Pic(X)=\mathbb Z[L]\) and \(L\)
  ample of even genus \(g=2k\).  For the Lazarsfeld--Mukai bundle \(E\), each
  nonzero section \(s\) determines a one-dimensional subspace
  \[
    \K_{k-1,1}\bigl(X,L,H^0(L\otimes I_{Z(s)})\bigr)\subseteq \K_{k-1,1}(X,L).
  \]
  The induced map
  \(\PP(H^0(E))\to \PP(\K_{k-1,1}(X,L))\) is the degree \(k-2\) Veronese
  embedding, hence induces
  \(\Sym^{k-2}H^0(X,E)\simeq \K_{k-1,1}(X,L)\).
\end{theorem}
\begin{proof}
  \uses{cor:kemeny_extremal_syzygy_iso, thm:kemeny_veronese_identification}
  The surjectivity theorem for
  \(H^1(B,\wedge^k\pi^*\cM\otimes p'^*L)\to H^1(B,\wedge^k\Gamma\otimes p'^*L)\)
  and the determinant computation for \(\Gamma\) identify the target cohomology
  with \(H^0(\PP,\mathcal O_{\PP}(k-2))^\vee\).  Dualizing gives an injection
  \(\Sym^{k-2}H^0(E)\hookrightarrow \K_{k-1,1}(X,L)\), and the even-genus
  vanishing gives equality of dimensions.  The fibre calculation in Section 2
  identifies the hyperplane quotient over \([t]\) with the secant Koszul
  subspace attached to \(Z(t)\), so the projective morphism is the Veronese map.
\end{proof}

\begin{corollary}[Geometric Syzygy Conjecture in even genus]
  \label{cor:kemeny_geometric_syzygy_even}
  \lean{MR4213770.GeometricSyzygyEven}
  \uses{thm:kemeny_even_geometric_syzygy_structure, lem:kemeny_grassmann_pencil_map, lem:kemeny_voisin_generation_observation, lem:kemeny_restrict_k3_syzygies_to_curve}
  % SOURCE: Kemeny, Introduction, Corollary geo-cor (read from references/kemeny-universal-secant-bundles-syzygies-canonical-curves-source/kemeny-universal-secant-bundles-syzygies-canonical-curves.tex)
  % SOURCE QUOTE: Let $C$ be a general curve of even genus $g=2k$. Then $\mathrm{K}_{k-1,1}(C,\omega_C)$ is generated by syzygies of the lowest possible rank $k$. More precisely, $\mathrm{K}_{k-1,1}(C,\omega_C)$ is generated by the rank $k$ syzygies $$\alpha \in \mathrm{K}_{k-1,1}(C,\omega_C,\mathrm{H}^0(\omega_C \otimes A^{-1})), \; \; A \in W^1_{k+1}(C).$$
  \textit{Source: Kemeny, Introduction, Corollary geo-cor.}
  If \(C\) is a general curve of even genus \(2k\), then
  \(\K_{k-1,1}(C,\omega_C)\) is generated by the rank \(k\) syzygies arising
  from the subspaces \(H^0(\omega_C\otimes A^{-1})\), for
  \(A\in W^1_{k+1}(C)\).
\end{corollary}
\begin{proof}
  \uses{thm:kemeny_even_geometric_syzygy_structure, lem:kemeny_grassmann_pencil_map, lem:kemeny_voisin_generation_observation, lem:kemeny_restrict_k3_syzygies_to_curve}
  For a general \(C\in |L|\), the finite map from
  \(\Gr_2(H^0(E))\) to \(|L|\) identifies its fibre over \(C\) with
  \(W^1_{k+1}(C)\).  The degree \(k-2\) Veronese image is nondegenerate on the
  union of the lines corresponding to these pencils, so the K3 space
  \(\K_{k-1,1}(X,L)\) is generated by the corresponding secant subspaces.
  Restriction to \(C\) identifies these subspaces with the rank \(k\) syzygy
  spaces in the statement.
\end{proof}

\begin{lemma}[Even K3 section transfer to a general curve]
  \label{lem:kemeny_even_k3_section_to_general_curve}
  \lean{MR4213770.EvenK3SectionToGeneralCurve}
  \uses{thm:kemeny_even_voisin}
  % SOURCE: Kemeny, Introduction, even-genus transfer sentence (read from references/kemeny-universal-secant-bundles-syzygies-canonical-curves-source/kemeny-universal-secant-bundles-syzygies-canonical-curves.tex)
  % SOURCE QUOTE: Voisin's Theorem states that $\mathrm{K}_{k,1}(X,L)=0$, \cite{V1}. This single vanishing suffices to prove Green's Conjecture for general canonical curves in even genus.
  \textit{Source: Kemeny, Introduction, even-genus transfer sentence.}
  The vanishing \(\K_{k,1}(X,L)=0\) for a Picard-rank-one K3 surface of even
  genus \(2k\) implies \(\K_{k,1}(C,\omega_C)=0\) for a general canonical curve
  \(C\) of genus \(2k\).
\end{lemma}
\begin{proof}
  \uses{thm:kemeny_even_voisin}
  Restrict the K3 Koszul complex to smooth hyperplane sections \(C\in |L|\);
  adjunction identifies \(L|_C\) with \(\omega_C\), and the standard Green
  restriction comparison identifies the relevant curve Koszul group with the
  specialization controlled by \(\K_{k,1}(X,L)\).  Since Kemeny's even K3
  vanishing supplies zero for this specialization and Koszul cohomology
  dimensions are upper semicontinuous in flat families of canonical curves, the
  same vanishing holds on the open general even-genus locus.
\end{proof}

\subsection{Voisin's Theorem in Even Genus}

\begin{definition}[Kernel bundle and Koszul cohomology reduction]
  \label{def:kemeny_kernel_bundle_reduction}
  \lean{MR4213770.KernelBundleReduction}
  \uses{def:peer_scheme_hmodule}
  % SOURCE: Kemeny, Section sec1, opening paragraph (read from references/kemeny-universal-secant-bundles-syzygies-canonical-curves-source/kemeny-universal-secant-bundles-syzygies-canonical-curves.tex)
  % SOURCE QUOTE: Let $M_L$ denote the bundle defined by the sequence $0 \to M_L \to \mathrm{H}^0(L)\otimes \mathcal{O}_X \to L \to 0.$ By the Kernel Bundle description of Koszul cohomology, see \cite{lazarsfeld-VBT} or \cite[\S 3]{ein-lazarsfeld-asymptotic}, to prove Voisin's Theorem it suffices to show $$\mathrm{H}^1(X, \bigwedge^{k+1} M_L )=0.$$
  \textit{Source: Kemeny, Section sec1, opening paragraph.}
  The kernel bundle \(M_L\) is defined by
  \(0\to M_L\to H^0(L)\otimes \cO_X\to L\to0\).  For the even-genus K3
  surface, the desired vanishing \(\K_{k,1}(X,L)=0\) follows from
  \(H^1(X,\wedge^{k+1}M_L)=0\).
\end{definition}

\begin{definition}[Universal zero locus]
  \label{def:kemeny_universal_zero_locus}
  \lean{MR4213770.UniversalZeroLocus}
  % SOURCE: Kemeny, Section sec1, universal construction (read from references/kemeny-universal-secant-bundles-syzygies-canonical-curves-source/kemeny-universal-secant-bundles-syzygies-canonical-curves.tex)
  % SOURCE QUOTE: Set $\PP:=\PP(\mathrm{H}^0(E))\simeq \PP^{k+1}$. Consider $X \times \PP$ with projections $p: X \times \PP \to X$, $q: X \times \PP \to \PP$. Define $\mathcal{Z} \seq X \times \PP$ as the locus $\left\{ (x,s) \; | \; s(x)=0 \right\}$.
  \textit{Source: Kemeny, Section sec1.}
  Let \(\PP=\PP(H^0(E))\), with projections \(p:X\times\PP\to X\) and
  \(q:X\times\PP\to\PP\).  The universal zero locus is
  \[
    \cZ=\{(x,s)\mid s(x)=0\}\subseteq X\times\PP .
  \]
  It is smooth over \(X\), and finite flat over \(\PP\) under the Picard-rank-one
  hypothesis.
\end{definition}

\begin{definition}[Universal secant bundles]
  \label{def:kemeny_universal_secant_bundles}
  \lean{MR4213770.UniversalSecantBundles}
  \uses{def:kemeny_universal_zero_locus}
  % SOURCE: Kemeny, Section sec1, definition after Lemma very-first-lem (read from references/kemeny-universal-secant-bundles-syzygies-canonical-curves-source/kemeny-universal-secant-bundles-syzygies-canonical-curves.tex)
  % SOURCE QUOTE: Let $\mathcal{S}$ be the vector bundle on $B$ defined by the exact sequence $$0 \to \mathcal{S} \to {q'}^* q'_* ({p'}^*L \otimes I_D) \to {p'}^*L \otimes I_D \to 0.$$ We call $\mathcal{S}$ and $\Gamma$ the \emph{universal secant bundles}. We have an exact sequence $$0 \to \mathcal{S} \to \pi^* \mathcal{M} \to \Gamma \to 0$$
  \textit{Source: Kemeny, Section sec1.}
  Let \(\pi:B\to X\times\PP\) be the blow-up of \(\cZ\), with exceptional
  divisor \(D\), and let \(p'=p\circ\pi\), \(q'=q\circ\pi\).  Define
  \[
    0\to\cS\to q'^*q'_*(p'^*L\otimes I_D)\to p'^*L\otimes I_D\to0
  \]
  and let \(\Gamma\) be the quotient in
  \(0\to\cS\to\pi^*\cM\to\Gamma\to0\).  The bundles \(\cS\) and \(\Gamma\) are
  Kemeny's universal secant bundles.
\end{definition}

\begin{lemma}[The cokernel bundle \(\mathcal W\)]
  \label{lem:kemeny_W_locally_free}
  \lean{MR4213770.WLocallyFree}
  \uses{def:kemeny_universal_zero_locus, def:peer_higher_direct_image, def:peer_affine_pushforward_base_change, def:peer_module_iso_locality}
  % SOURCE: Kemeny, Section sec1, Lemma very-first-lem (read from references/kemeny-universal-secant-bundles-syzygies-canonical-curves-source/kemeny-universal-secant-bundles-syzygies-canonical-curves.tex)
  % SOURCE QUOTE: The sheaf $\mathcal{W}$ is locally free of rank $k$.
  \textit{Source: Kemeny, Lemma very-first-lem.}
  The cokernel
  \[
    \mathcal W=\operatorname{Coker}\bigl(q'_*(p'^*L\otimes I_D)\to q'_*p'^*L\bigr)
  \]
  is locally free of rank \(k\).
\end{lemma}
\begin{proof}
  \uses{def:peer_higher_direct_image, def:peer_affine_pushforward_base_change, def:peer_module_iso_locality}
  Apply \(Rq_*\) to the universal exact sequence.  The fibres of
  \(q_*(p^*L|_{\cZ})\) have dimension \(k+1\), while
  \(R^1q_*(p^*L\otimes I_{\cZ})\) has fibre \(H^2(\cO_X)\), hence rank one.
  Grauert's theorem gives local freeness.  The pushforward exact sequence is
  checked after the standard base-change reductions, and the resulting
  sheaf-map identification is local on stalks; the exact sequence then shows
  that the kernel \(\mathcal W\) has rank \(k\).
\end{proof}

\begin{lemma}[First symmetric-power computation]
  \label{lem:kemeny_first_iso}
  \lean{MR4213770.FirstIso}
  \uses{def:kemeny_universal_secant_bundles}
  % SOURCE: Kemeny, Section sec1, Lemma first-iso (read from references/kemeny-universal-secant-bundles-syzygies-canonical-curves-source/kemeny-universal-secant-bundles-syzygies-canonical-curves.tex)
  % SOURCE QUOTE: We have $\mathrm{H}^{k+1}(X \times \PP,\mathrm{Sym}^{k+1} \mathcal{G})=0$ as well as a natural isomorphism $$\mathrm{H}^k(X \times \PP,\mathrm{Sym}^{k+1} \mathcal{G})\simeq \mathrm{Sym}^k \mathrm{H}^0(E).$$
  \textit{Source: Kemeny, Lemma first-iso.}
  For \(\cG=q^*q_*(p^*L\otimes I_{\cZ})\), one has
  \(H^{k+1}(X\times\PP,\Sym^{k+1}\cG)=0\) and a natural isomorphism
  \(H^k(X\times\PP,\Sym^{k+1}\cG)\simeq \Sym^kH^0(E)\).
\end{lemma}
\begin{proof}
  Resolve \(\Sym^{k+1}\cG\) by the two-term symmetric-power sequence induced
  from \(0\to q^*\cO_{\PP}(-2)\to H^0(E)\otimes q^*\cO_{\PP}(-1)\to\cG\to0\).
  The cohomology of negative twists on \(\PP^{k+1}\), Kunneth, and
  \(H^1(\cO_X)=0\) identify the only nonzero contribution in degree \(k\) with
  \(\Sym^kH^0(E)\) and force the degree \(k+1\) group to vanish.
\end{proof}

\begin{lemma}[Twisted symmetric-power comparison]
  \label{lem:kemeny_twisted_symmetric_comparison}
  \lean{MR4213770.TwistedSymmetricComparison}
  \uses{lem:kemeny_first_iso}
  % SOURCE: Kemeny, Section sec1, Lemmas first-twist-lem and second-twist-lem (read from references/kemeny-universal-secant-bundles-syzygies-canonical-curves-source/kemeny-universal-secant-bundles-syzygies-canonical-curves.tex)
  % SOURCE QUOTE: The evaluation morphism $\mathcal{G} \twoheadrightarrow p^*L \otimes I_{\mathcal{Z}}$ induces an isomorphism
  \textit{Source: Kemeny, Lemmas first-twist-lem and second-twist-lem.}
  The group \(H^k(\Sym^k\cG\otimes p^*L\otimes I_{\cZ})\) is naturally
  \(\Sym^kH^0(E)\), and the evaluation map \(\cG\twoheadrightarrow
  p^*L\otimes I_{\cZ}\) induces an isomorphism
  \[
    H^k(\Sym^k\cG\otimes\cG)\simeq
    H^k(\Sym^k\cG\otimes p^*L\otimes I_{\cZ}).
  \]
\end{lemma}
\begin{proof}
  \uses{lem:kemeny_first_iso}
  Tensor the symmetric-power resolution of \(\cG\) with \(p^*L\otimes I_{\cZ}\).
  The universal section exact sequence identifies the relevant boundary map with
  the identity of \(H^0(E)\), so the only surviving degree \(k\) contribution is
  \(\Sym^kH^0(E)\).  Repeating the same computation with \(\cG\) in place of
  \(p^*L\otimes I_{\cZ}\) gives the same vector space, and the evaluation map is
  the identity under these identifications.
\end{proof}

\begin{proposition}[Automatic injectivity]
  \label{prop:kemeny_auto_injectivity}
  \lean{MR4213770.AutoInjectivity}
  \uses{lem:kemeny_first_iso, lem:kemeny_twisted_symmetric_comparison}
  % SOURCE: Kemeny, Section sec1, Proposition auto-injectivity (read from references/kemeny-universal-secant-bundles-syzygies-canonical-curves-source/kemeny-universal-secant-bundles-syzygies-canonical-curves.tex)
  % SOURCE QUOTE: We have a natural isomorphism $\mathrm{H}^k(\mathrm{Sym}^{k+1} \mathcal{G}) \xrightarrow{\sim} \mathrm{H}^k(\mathrm{Sym}^{k} \mathcal{G} \otimes p^*L \otimes I_{\mathcal{Z}}).$
  \textit{Source: Kemeny, Proposition auto-injectivity.}
  The natural map
  \(H^k(\Sym^{k+1}\cG)\to H^k(\Sym^k\cG\otimes p^*L\otimes I_{\cZ})\) is an
  isomorphism.
\end{proposition}
\begin{proof}
  \uses{lem:kemeny_first_iso, lem:kemeny_twisted_symmetric_comparison}
  Under the two preceding identifications, it is enough to prove injectivity of
  \(H^k(\Sym^{k+1}\cG)\to H^k(\Sym^k\cG\otimes\cG)\).  For any vector bundle
  \(\mathcal F\), the composite
  \(\Sym^i\mathcal F\to \Sym^{i-1}\mathcal F\otimes\mathcal F\to\Sym^i\mathcal F\)
  is multiplication by \(i\), hence is injective over \(\C\).
\end{proof}

\begin{theorem}[Even-genus secant-bundle vanishing]
  \label{thm:kemeny_even_secant_vanishing}
  \lean{MR4213770.EvenSecantVanishing}
  \uses{def:kemeny_universal_secant_bundles, prop:kemeny_auto_injectivity, def:peer_affine_pushforward_base_change}
  % SOURCE: Kemeny, Section sec1, Theorem main-thm (read from references/kemeny-universal-secant-bundles-syzygies-canonical-curves-source/kemeny-universal-secant-bundles-syzygies-canonical-curves.tex)
  % SOURCE QUOTE: We have $\mathrm{H}^i(B, \bigwedge^{k+1-i} \pi^* \mathcal{M} \otimes \mathrm{Sym}^i \mathcal{S})=0$ for $1 \leq i \leq k+1$.
  \textit{Source: Kemeny, Theorem main-thm.}
  For \(1\le i\le k+1\),
  \[
    H^i\bigl(B,\wedge^{k+1-i}\pi^*\cM\otimes\Sym^i\cS\bigr)=0 .
  \]
\end{theorem}
\begin{proof}
  \uses{prop:kemeny_auto_injectivity, def:peer_affine_pushforward_base_change}
  Use the symmetric-power sequence for the quotient defining \(\cS\).  After
  applying the projection formula and the blow-up identities
  \(\pi_*\cO_B=\cO_{X\times\PP}\), \(\pi_*I_D=I_{\cZ}\), and vanishing of higher
  direct images, the top case is exactly the automatic-injectivity isomorphism.
  The remaining cases reduce to Kunneth vanishings for exterior powers of
  \(M_L\) tensored with negative line bundles on \(\PP\).
\end{proof}

\begin{theorem}[Voisin's theorem in even genus]
  \label{thm:kemeny_even_voisin}
  \lean{MR4213770.EvenVoisin}
  \uses{def:kemeny_kernel_bundle_reduction, thm:kemeny_even_secant_vanishing}
  % SOURCE: Kemeny, Section sec1 and Introduction (read from references/kemeny-universal-secant-bundles-syzygies-canonical-curves-source/kemeny-universal-secant-bundles-syzygies-canonical-curves.tex)
  % SOURCE QUOTE: We now complete the proof that $\mathrm{K}_{k,1}(X,L)=0$.
  \textit{Source: Kemeny, Section sec1.}
  For a Picard-rank-one complex K3 surface \((X,L)\) of even genus \(2k\),
  \(\K_{k,1}(X,L)=0\).
\end{theorem}
\begin{proof}
  \uses{def:kemeny_kernel_bundle_reduction, thm:kemeny_even_secant_vanishing}
  The kernel-bundle description reduces \(\K_{k,1}(X,L)=0\) to
  \(H^1(X,\wedge^{k+1}M_L)=0\).  Pulling to \(B\) and using the exact sequence
  \(0\to\cS\to\pi^*\cM\to\Gamma\to0\), the exterior-power filtration reduces
  this vanishing to the family of vanishing statements in
  Theorem~\ref{thm:kemeny_even_secant_vanishing}.
\end{proof}

\subsection{The Geometric Syzygy Conjecture}

\begin{theorem}[Surjectivity for the extremal syzygy map]
  \label{thm:kemeny_real_thm}
  \lean{MR4213770.ExtremalMapSurjective}
  \uses{def:kemeny_universal_secant_bundles}
  % SOURCE: Kemeny, Section sec2, Theorem real-thm (read from references/kemeny-universal-secant-bundles-syzygies-canonical-curves-source/kemeny-universal-secant-bundles-syzygies-canonical-curves.tex)
  % SOURCE QUOTE: The natural morphism $\psi: \mathrm{H}^1(B, \wedge^k \pi^* \mathcal{M} \otimes {p'}^*L) \to \mathrm{H}^1(B, \wedge^k \Gamma \otimes {p'}^*L)$ is surjective.
  \textit{Source: Kemeny, Theorem real-thm.}
  The natural map
  \(H^1(B,\wedge^k\pi^*\cM\otimes p'^*L)\to
    H^1(B,\wedge^k\Gamma\otimes p'^*L)\)
  is surjective.
\end{theorem}
\begin{proof}
  \uses{def:kemeny_universal_secant_bundles}
  The exterior-power sequence for \(0\to\cS\to\pi^*\cM\to\Gamma\to0\) reduces
  surjectivity to the vanishing of
  \(H^{1+i}(\Sym^i\cS\otimes\wedge^{k-i}\pi^*\cM\otimes p'^*L)\).
  The symmetric-power sequence for \(\cS\), the projection formula, and the
  universal zero-locus exact sequence reduce these groups to Kunneth vanishings
  on \(X\times\PP\).
\end{proof}

\begin{lemma}[Determinant of \(\Gamma\) in even genus]
  \label{lem:kemeny_even_det_gamma}
  \lean{MR4213770.EvenDetGamma}
  \uses{lem:kemeny_W_locally_free}
  % SOURCE: Kemeny, Section sec2, determinant lemma after Theorem real-thm (read from references/kemeny-universal-secant-bundles-syzygies-canonical-curves-source/kemeny-universal-secant-bundles-syzygies-canonical-curves.tex)
  % SOURCE QUOTE: With notation as in Section \ref{sec1}, we have $\wedge^k \Gamma \simeq I_D \otimes {q'}^* \mathcal{O}_{\PP}(k)$.
  \textit{Source: Kemeny, Section sec2, determinant lemma.}
  With the notation of the even-genus construction,
  \(\wedge^k\Gamma\simeq I_D\otimes q'^*\cO_{\PP}(k)\).
\end{lemma}
\begin{proof}
  \uses{lem:kemeny_W_locally_free}
  Take determinants in
  \(0\to\Gamma\to q'^*\mathcal W\to p'^*L|_D\to0\).  The determinant of
  \(\mathcal W\) is computed from
  \(0\to q_*(p^*L\otimes I_{\cZ})\to q_*p^*L\to\mathcal W\to0\), and the
  two-term resolution of \(q_*(p^*L\otimes I_{\cZ})\) gives
  \(\det\mathcal W\simeq \cO_{\PP}(k)\).
\end{proof}

\begin{corollary}[Extremal syzygy vector-space isomorphism]
  \label{cor:kemeny_extremal_syzygy_iso}
  \lean{MR4213770.ExtremalSyzygyIso}
  \uses{thm:kemeny_real_thm, lem:kemeny_even_det_gamma, thm:kemeny_even_voisin}
  % SOURCE: Kemeny, Section sec2, Corollary main-cor-nat-iso (read from references/kemeny-universal-secant-bundles-syzygies-canonical-curves-source/kemeny-universal-secant-bundles-syzygies-canonical-curves.tex)
  % SOURCE QUOTE: The morphism $\psi: \mathrm{H}^1(\wedge^k \pi^* \mathcal{M} \otimes {p'}^*L) \to \mathrm{H}^1(\wedge^k \Gamma \otimes {p'}^*L)$ is an isomorphism and induces a natural isomorphism $\mathrm{K}_{k-1,1}(X,L) \simeq \mathrm{Sym}^{k-2}\mathrm{H}^0(X,E)  $.
  \textit{Source: Kemeny, Corollary main-cor-nat-iso.}
  The map \(\psi\) is an isomorphism and induces a natural isomorphism
  \(\K_{k-1,1}(X,L)\simeq \Sym^{k-2}H^0(X,E)\).
\end{corollary}
\begin{proof}
  \uses{thm:kemeny_real_thm, lem:kemeny_even_det_gamma, thm:kemeny_even_voisin}
  The determinant computation rewrites the target of \(\psi\) as the cohomology
  of \(I_D\otimes q'^*\cO_{\PP}(k)\otimes p'^*L\), which is
  \(H^0(\cO_{\PP}(k-2))^\vee\).  Surjectivity gives the dual injection
  \(\Sym^{k-2}H^0(E)\hookrightarrow \K_{k-1,1}(X,L)\); the even-genus vanishing
  supplies the dimension equality.
\end{proof}

\begin{lemma}[Image of \(\wedge^k\Gamma'\)]
  \label{lem:kemeny_gamma_prime_image}
  \lean{MR4213770.GammaPrimeImage}
  \uses{def:kemeny_universal_secant_bundles, lem:kemeny_W_locally_free}
  % SOURCE: Kemeny, Section sec2, Proof of Theorem geo-thm-intro (read from references/kemeny-universal-secant-bundles-syzygies-canonical-curves-source/kemeny-universal-secant-bundles-syzygies-canonical-curves.tex)
  % SOURCE QUOTE: We have the exact sequence $0 \to \Gamma' \to q^* \mathcal{W} \to p^*L_{|_{\mathcal{Z}}} \to 0,$ where we have set $\Gamma':=\pi_* \Gamma$. The image of the natural map $\wedge^k \Gamma' \to \wedge^k q^* \mathcal{W}  \simeq \mathcal{O}_{\PP}(k)$ is $I_{\mathcal{Z}} \otimes q^*\mathcal{O}_{\PP}(k)$ from \cite[Cor.\ 3.7]{ein-lazarsfeld-asymptotic}.
  \textit{Source: Kemeny, Proof of Theorem geo-thm-intro.}
  With \(\Gamma'=\pi_*\Gamma\), there is an exact sequence
  \(0\to\Gamma'\to q^*\mathcal W\to p^*L|_{\cZ}\to0\).  The image of
  \(\wedge^k\Gamma'\to\wedge^kq^*\mathcal W\simeq\cO_{\PP}(k)\) is
  \(I_{\cZ}\otimes q^*\cO_{\PP}(k)\), and the induced map is the pushforward of
  \(\wedge^k\Gamma\to I_D\otimes q'^*\cO_{\PP}(k)\).
\end{lemma}
\begin{proof}
  \uses{def:kemeny_universal_secant_bundles, lem:kemeny_W_locally_free}
  Push forward the quotient sequence defining \(\Gamma\) and compare its
  determinant with the determinant of \(q^*\mathcal W\).  The standard secant
  image computation for the zero locus \(\cZ\) identifies the resulting ideal
  image with \(I_{\cZ}\otimes q^*\cO_{\PP}(k)\).
\end{proof}

\begin{lemma}[The \(R^1q_*\)-line bundle]
  \label{lem:kemeny_R1q_line_bundle}
  \lean{MR4213770.R1qLineBundle}
  \uses{def:kemeny_universal_zero_locus, lem:kemeny_W_locally_free, def:peer_higher_direct_image, def:peer_affine_pushforward_base_change}
  % SOURCE: Kemeny, Section sec2, Proof of Theorem geo-thm-intro (read from references/kemeny-universal-secant-bundles-syzygies-canonical-curves-source/kemeny-universal-secant-bundles-syzygies-canonical-curves.tex)
  % SOURCE QUOTE: From $0 \to \mathcal{O}_X \boxtimes \mathcal{O}_{\PP}(-2) \to E \boxtimes \mathcal{O}_{\PP}(-1) \to p^*L \otimes  I_{\mathcal{Z}} \to 0,$ we have $$R^1 q_*((L \boxtimes \mathcal{O}_{\PP}(k))\otimes I_{\mathcal{Z}}) \simeq R^2q_*\mathcal{O}_{X \times \PP} \otimes \mathcal{O}_{\PP}(k-2)\simeq \mathcal{O}_{\PP}(k-2),$$
  \textit{Source: Kemeny, Proof of Theorem geo-thm-intro.}
  The higher direct image
  \(R^1q_*((L\boxtimes\cO_{\PP}(k))\otimes I_{\cZ})\) is canonically
  isomorphic to \(\cO_{\PP}(k-2)\).
\end{lemma}
\begin{proof}
  \uses{def:kemeny_universal_zero_locus, lem:kemeny_W_locally_free, def:peer_higher_direct_image, def:peer_affine_pushforward_base_change}
  Twist the universal zero-locus sequence by \(q^*\cO_{\PP}(k)\) and apply
  \(Rq_*\).  The middle term has no relevant higher image, and relative duality
  gives \(R^2q_*\cO_{X\times\PP}\simeq\cO_{\PP}\), producing the line bundle
  \(\cO_{\PP}(k-2)\).
\end{proof}

\begin{lemma}[Global-section description of \(\psi\)]
  \label{lem:kemeny_psi_global_sections}
  \lean{MR4213770.PsiGlobalSections}
  \uses{cor:kemeny_extremal_syzygy_iso, lem:kemeny_gamma_prime_image, lem:kemeny_R1q_line_bundle, def:peer_higher_direct_image}
  % SOURCE: Kemeny, Section sec2, Proof of Theorem geo-thm-intro (read from references/kemeny-universal-secant-bundles-syzygies-canonical-curves-source/kemeny-universal-secant-bundles-syzygies-canonical-curves.tex)
  % SOURCE QUOTE: The map $\psi$ is thus naturally identified with the global section map applied to the morphism $$\mathrm{K}_{k-1,1}(X,L)^{\vee} \otimes \mathcal{O}_{\PP} \simeq R^1q_*\wedge^k\mathcal{M} \otimes p^*L \to R^1 q_*((L \boxtimes \mathcal{O}_{\PP}(k))\otimes I_{\mathcal{Z}}),$$
  \textit{Source: Kemeny, Proof of Theorem geo-thm-intro.}
  Under Leray and the preceding image computation, \(\psi\) is the map on global
  sections induced by
  \[
    \K_{k-1,1}(X,L)^\vee\otimes\cO_{\PP}\to\cO_{\PP}(k-2).
  \]
\end{lemma}
\begin{proof}
  \uses{cor:kemeny_extremal_syzygy_iso, lem:kemeny_gamma_prime_image, lem:kemeny_R1q_line_bundle, def:peer_higher_direct_image}
  Leray rewrites the source cohomology as global sections of
  \(R^1q_*(\wedge^k\mathcal M\otimes p^*L)\).  The map
  \(\wedge^k\Gamma'\to I_{\cZ}\otimes q^*\cO_{\PP}(k)\), followed by the
  \(R^1q_*\)-identification, gives the displayed vector-bundle morphism.
\end{proof}

\begin{lemma}[Fibre of the secant map]
  \label{lem:kemeny_fibre_secant_map}
  \lean{MR4213770.FibreSecantMap}
  \uses{lem:kemeny_psi_global_sections, def:kemeny_universal_secant_bundles}
  % SOURCE: Kemeny, Section sec2, Proof of Theorem geo-thm-intro (read from references/kemeny-universal-secant-bundles-syzygies-canonical-curves-source/kemeny-universal-secant-bundles-syzygies-canonical-curves.tex)
  % SOURCE QUOTE: The fiber of the morphism (\ref{bundle-ind-vero}) over $[t]$ is identified with the natural composition $$\mathrm{H}^1(X,\wedge^k M_L(L)) \to \mathrm{H}^1(X,\wedge^k \Gamma'_t(L)) \to \mathrm{H}^1(X,L \otimes I_{Z(t)}),$$
  \textit{Source: Kemeny, Proof of Theorem geo-thm-intro.}
  For nonzero \(t\in H^0(E)\), the fibre over \([t]\) of the morphism
  \(\K_{k-1,1}(X,L)^\vee\otimes\cO_{\PP}\to\cO_{\PP}(k-2)\) is the map through
  the blow-up \(B_t\) of \(Z(t)\).  Under Serre duality and the kernel-bundle
  description, its dual image is precisely
  \[
    \K_{k-1,1}\bigl(X,L,H^0(L\otimes I_{Z(t)})\bigr)
    \subseteq \K_{k-1,1}(X,L).
  \]
\end{lemma}
\begin{proof}
  \uses{lem:kemeny_psi_global_sections, def:kemeny_universal_secant_bundles}
  Evaluate the global vector-bundle morphism at \([t]\) and identify the fibre
  with the cohomology map on the individual secant blow-up \(B_t\).  The
  determinant identifications convert this to the Serre dual inclusion
  \(\wedge^k\mathcal S_t(D_t)\hookrightarrow\wedge^k\pi_t^*M_L(D_t)\), which is
  the inclusion of the secant Koszul subspace.
\end{proof}

\begin{theorem}[Veronese identification of the secant syzygy map]
  \label{thm:kemeny_veronese_identification}
  \lean{MR4213770.VeroneseIdentification}
  \uses{cor:kemeny_extremal_syzygy_iso, lem:kemeny_psi_global_sections, lem:kemeny_fibre_secant_map}
  % SOURCE: Kemeny, Section sec2, Proof of Theorem geo-thm-intro (read from references/kemeny-universal-secant-bundles-syzygies-canonical-curves-source/kemeny-universal-secant-bundles-syzygies-canonical-curves.tex)
  % SOURCE QUOTE: is surjective on global sections and $\mathcal{O}_{\PP}(k-2)$ is globally generated, it must be surjective as a morphism of sheaves.
  \textit{Source: Kemeny, Proof of Theorem geo-thm-intro.}
  The morphism
  \(\K_{k-1,1}(X,L)^\vee\otimes\cO_{\PP}\to\cO_{\PP}(k-2)\) induced by \(\psi\)
  is a quotient bundle.  The associated projective map
  \(\PP(H^0(E))\to\PP(\K_{k-1,1}(X,L))\) is the degree \(k-2\) Veronese map, and
  its point over \([t]\) is the secant Koszul line associated to \(Z(t)\).
\end{theorem}
\begin{proof}
  \uses{cor:kemeny_extremal_syzygy_iso, lem:kemeny_psi_global_sections, lem:kemeny_fibre_secant_map}
  The global-section description identifies the map as a morphism onto
  \(\cO_{\PP}(k-2)\).  The extremal syzygy isomorphism makes it surjective on
  global sections; since \(\cO_{\PP}(k-2)\) is globally generated, the morphism is
  surjective as a sheaf map.  Its projectivization is therefore the Veronese
  quotient, and the fibre lemma identifies its points with the desired secant
  Koszul subspaces.
\end{proof}

\begin{lemma}[Grassmannian map to \(|L|\)]
  \label{lem:kemeny_grassmann_pencil_map}
  \lean{MR4213770.GrassmannPencilMap}
  % SOURCE: Kemeny, Section sec2, Proof of Corollary geo-cor (read from references/kemeny-universal-secant-bundles-syzygies-canonical-curves-source/kemeny-universal-secant-bundles-syzygies-canonical-curves.tex)
  % SOURCE QUOTE: We have a surjective, finite morphism $$d: \mathrm{Gr}_2(\mathrm{H}^0(E)) \to |L|$$ which takes a two dimensional subspace $W \seq \mathrm{H}^0(E)$ to the degeneracy locus of the evaluation morphism $W \otimes \mathcal{O}_X \xrightarrow{ev} E$, \cite{V1}.
  \textit{Source: Kemeny, Proof of Corollary geo-cor.}
  There is a finite surjective morphism
  \(d:\Gr_2(H^0(E))\to |L|\) sending \(W\subset H^0(E)\) to the degeneracy
  divisor of \(W\otimes\cO_X\to E\).  For a general \(C\in|L|\), the fibre
  \(d^{-1}(C)\) is naturally identified with the reduced Brill--Noether locus
  \(W^1_{k+1}(C)\).
\end{lemma}
\begin{proof}
  The evaluation degeneracy construction gives the morphism to \(|L|\).  For a
  general curve \(C\), the cokernel of evaluation is \(\omega_C\otimes A^{-1}\),
  and the Brill--Noether comparison identifies \(W\) with \(H^0(A)^\vee\) for
  \(A\in W^1_{k+1}(C)\).
\end{proof}

\begin{lemma}[Voisin generation observation]
  \label{lem:kemeny_voisin_generation_observation}
  \lean{MR4213770.VoisinGenerationObservation}
  \uses{lem:kemeny_grassmann_pencil_map}
  % SOURCE: Kemeny, Section sec2, Proof of Corollary geo-cor (read from references/kemeny-universal-secant-bundles-syzygies-canonical-curves-source/kemeny-universal-secant-bundles-syzygies-canonical-curves.tex)
  % SOURCE QUOTE: By \cite[Proof of Prop.\ 7]{V1}, the spaces $\rm{Sym}^{k-2}\mathrm{H}^0(A)^{\vee}$, $A \in W^1_{k+1}(C)$ generate $\rm{Sym}^{k-2}\mathrm{H}^0(E)$.
  \textit{Source: Kemeny, Proof of Corollary geo-cor.}
  For general \(C\in |L|\), the subspaces
  \(\Sym^{k-2}H^0(A)^\vee\), as \(A\) ranges over \(W^1_{k+1}(C)\), generate
  \(\Sym^{k-2}H^0(E)\).
\end{lemma}
\begin{proof}
  \uses{lem:kemeny_grassmann_pencil_map}
  Under the fibre identification \(d^{-1}(C)\simeq W^1_{k+1}(C)\), the pencils
  give the lines \(H^0(A)^\vee\subset H^0(E)\).  Voisin's generation statement
  says that their \((k-2)\)-nd symmetric powers span the full symmetric power.
\end{proof}

\begin{lemma}[Restriction of K3 secant syzygies]
  \label{lem:kemeny_restrict_k3_syzygies_to_curve}
  \lean{MR4213770.RestrictK3SyzygiesToCurve}
  \uses{thm:kemeny_veronese_identification, lem:kemeny_grassmann_pencil_map}
  % SOURCE: Kemeny, Section sec2, Proof of Corollary geo-cor (read from references/kemeny-universal-secant-bundles-syzygies-canonical-curves-source/kemeny-universal-secant-bundles-syzygies-canonical-curves.tex)
  % SOURCE QUOTE: After restricting to $C$, such spaces are identified with subspaces of the form $$\mathrm{K}_{k-1,1}(C,\omega_C, \mathrm{H}^0(\omega_C \otimes A^{\vee})) \seq \mathrm{K}_{k-1,1}(C,\omega_C),\; \; A \in W^1_{k+1}(C)$$ under the isomorphism $\mathrm{K}_{k-1,1}(X,L) \simeq \mathrm{K}_{k-1,1}(C, \omega_C)$.
  \textit{Source: Kemeny, Proof of Corollary geo-cor.}
  Under restriction
  \(\K_{k-1,1}(X,L)\simeq \K_{k-1,1}(C,\omega_C)\), the K3 secant subspaces
  attached to sections corresponding to \(A\in W^1_{k+1}(C)\) become exactly
  \[
    \K_{k-1,1}\bigl(C,\omega_C,H^0(\omega_C\otimes A^{-1})\bigr)
    \subseteq \K_{k-1,1}(C,\omega_C).
  \]
\end{lemma}
\begin{proof}
  \uses{thm:kemeny_veronese_identification, lem:kemeny_grassmann_pencil_map}
  The section \(s\) corresponding to the pencil \(A\) has zero scheme whose
  secant subspace is identified by the Veronese fibre calculation.  Restricting
  the K3 Koszul complex to the hyperplane section \(C\) sends
  \(H^0(L\otimes I_{Z(s)})\) to \(H^0(\omega_C\otimes A^{-1})\), giving the
  stated curve subspace.
\end{proof}

\subsection{Voisin's Theorem in Odd Genus}

\begin{lemma}[Picard-rank-two K3 input]
  \label{lem:kemeny_BNP_gen}
  \lean{MR4213770.BNPGen}
  % SOURCE: Kemeny, Section odd-genus, Lemma BNP-gen (read from references/kemeny-universal-secant-bundles-syzygies-canonical-curves-source/kemeny-universal-secant-bundles-syzygies-canonical-curves.tex)
  % SOURCE QUOTE: Set $L_n:=L'-n\Delta$ for $0 \leq n \leq 3$. Then $L_n$ is base-point free and is further ample for $0<n \leq 2$. Furthermore, the linear system $|L_n|$ cannot be written as a sum of two pencils. In particular, smooth curves $C \in |L_n|$ are Brill--Noether general. If $C \in |L_n|$ is general then $C$ is, in addition, Petri general.
  \textit{Source: Kemeny, Lemma BNP-gen.}
  For the Picard-rank-two K3 surface used in the odd-genus reduction, the line
  bundles \(L_n=L'-n\Delta\), \(0\le n\le3\), are base-point free, \(L_n\) is
  ample for \(0<n\le2\), \(|L_n|\) is not a sum of two pencils, and general
  smooth curves in \(|L_n|\) are Brill--Noether and Petri general.
\end{lemma}
\begin{proof}
  Riemann--Roch gives nonempty systems.  A negative intersection rational base
  component would force an impossible decomposition in the rank-two Picard
  lattice, so \(L_n\) is nef; Mayer's criterion gives base-point freeness and
  the same lattice calculation gives ampleness for \(0<n\le2\).  If
  \(L_n=A_1+A_2\) with both \(h^0(A_i)\ge2\), expressing the \(A_i\) in the
  basis \(L',\Delta\) forces one summand to have only one section.
\end{proof}

\begin{lemma}[Projection criterion]
  \label{lem:kemeny_projection_criterion}
  \lean{MR4213770.ProjectionCriterion}
  \uses{lem:kemeny_BNP_gen}
  % SOURCE: Kemeny, Section odd-genus, Lemma proj-crit (read from references/kemeny-universal-secant-bundles-syzygies-canonical-curves-source/kemeny-universal-secant-bundles-syzygies-canonical-curves.tex)
  % SOURCE QUOTE: Suppose the projection map $pr_k$ is injective. Then $\mathrm{H}^1(\wedge^{k+1}M_L)=0$.
  \textit{Source: Kemeny, Lemma proj-crit.}
  If the projection map
  \(pr_k:H^1(\wedge^{k+1}M_{L'})\to H^1(\wedge^kM_L(-\Delta))\) is injective,
  then \(H^1(\wedge^{k+1}M_L)=0\).
\end{lemma}
\begin{proof}
  Take cohomology of
  \(0\to\wedge^{k+1}M_L\to\wedge^{k+1}M_{L'}\to\wedge^kM_L(-\Delta)\to0\).
  The group \(H^0(\wedge^kM_L(-\Delta))\) vanishes, so injectivity of \(pr_k\)
  kills the connecting obstruction and gives the desired \(H^1\)-vanishing.
\end{proof}

\begin{lemma}[Even-style result on the nodal K3 surface]
  \label{lem:kemeny_green_pic2}
  \lean{MR4213770.GreenPicTwo}
  \uses{thm:kemeny_even_voisin, cor:kemeny_extremal_syzygy_iso}
  % SOURCE: Kemeny, Section odd-genus, Lemma green-pic2 (read from references/kemeny-universal-secant-bundles-syzygies-canonical-curves-source/kemeny-universal-secant-bundles-syzygies-canonical-curves.tex)
  % SOURCE QUOTE: We have $\mathrm{K}_{k+1,1}(X,L')=0$. Further, there is a natural isomorphism $$\mathrm{Sym}^{k-1}\mathrm{H}^0(X,E) \simeq \mathrm{K}_{k,1}(X,L').$$
  \textit{Source: Kemeny, Lemma green-pic2.}
  For the odd-genus Picard-rank-two surface,
  \(\K_{k+1,1}(X,L')=0\) and
  \(\Sym^{k-1}H^0(X,E)\simeq \K_{k,1}(X,L')\).
  Here \(L=L'-\Delta\), so the blueprint's \(\K_{k,1}(X,L')\) is the paper's
  \(\K_{k,1}(X,L+\Delta)\).
\end{lemma}
\begin{proof}
  \uses{thm:kemeny_even_voisin, cor:kemeny_extremal_syzygy_iso}
  Contract \(\Delta\) to a rational double point on \(\hat X\).  The universal
  zero-locus construction remains finite flat because \(\operatorname{Cl}(\hat
  X)=\mathbb Z[\hat L]\), so the even-genus cohomological proof applies on
  \(\hat X\).  Rational singularities identify the section rings of \(L'\) and
  \(\hat L\), transferring the Koszul statements back to \(X\).
\end{proof}

\begin{lemma}[Nonvanishing on decomposable elements]
  \label{lem:kemeny_decomposable}
  \lean{MR4213770.DecomposableNonzero}
  \uses{lem:kemeny_BNP_gen}
  % SOURCE: Kemeny, Section odd-genus, Lemma decomposable (read from references/kemeny-universal-secant-bundles-syzygies-canonical-curves-source/kemeny-universal-secant-bundles-syzygies-canonical-curves.tex)
  % SOURCE QUOTE: The natural map $d: \wedge^2 \mathrm{H}^0(E) \to \mathrm{H}^0(\wedge^2 E)=\mathrm{H}^0(L')$ does not vanish on decomposable elements.
  \textit{Source: Kemeny, Lemma decomposable.}
  The determinant map
  \(\wedge^2H^0(E)\to H^0(\wedge^2E)=H^0(L')\) is nonzero on every nonzero
  decomposable element.
\end{lemma}
\begin{proof}
  \uses{lem:kemeny_BNP_gen}
  If a decomposable wedge vanished, the two sections would generate a rank-one
  subsheaf with at least two sections.  After removing torsion from the quotient,
  both determinant factors have at least two sections, so \(L'\) would split as a
  sum of two pencils, contradicting Lemma~\ref{lem:kemeny_BNP_gen}.
\end{proof}

\begin{lemma}[Twisting \(E\) by the nodal curve gives \(F\)]
  \label{lem:kemeny_E_minus_delta_is_F}
  \lean{MR4213770.EMinusDeltaIsF}
  \uses{lem:kemeny_BNP_gen}
  % SOURCE: Kemeny, Section odd-genus, stability argument before the definition of \(N\) (read from references/kemeny-universal-secant-bundles-syzygies-canonical-curves-source/kemeny-universal-secant-bundles-syzygies-canonical-curves.tex)
  % SOURCE QUOTE: The rank two bundle $E(-\Delta)$ is simple with $\det E(-\Delta)=L-\Delta$ and $\chi(E(-\Delta))=k+1$. We claim that $E(-\Delta)$ is stable with respect to $L-\Delta$. Since $F$ is the unique stable bundle with these invariants, $E(-\Delta) \simeq F$.
  \textit{Source: Kemeny, Section odd-genus.}
  The rank-two Lazarsfeld--Mukai bundle satisfies \(E(-\Delta)\simeq F\), where
  \(F\) is the unique stable bundle with determinant \(L-\Delta\) and Euler
  characteristic \(k+1\).
\end{lemma}
\begin{proof}
  \uses{lem:kemeny_BNP_gen}
  The bundle \(E(-\Delta)\) has the same numerical invariants as \(F\).  If it
  were unstable, a destabilizing rank-one filtration would give two line
  bundles with at least two sections whose sum is \(L\), contradicting the
  no-sum-of-pencils property in the rank-two Picard lattice.
\end{proof}

\begin{definition}[Odd-genus kernel bundle \(N\)]
  \label{def:kemeny_odd_kernel_N}
  \lean{MR4213770.OddKernelN}
  \uses{lem:kemeny_E_minus_delta_is_F}
  % SOURCE: Kemeny, Section odd-genus, definition of \(N\) and \(r\) (read from references/kemeny-universal-secant-bundles-syzygies-canonical-curves-source/kemeny-universal-secant-bundles-syzygies-canonical-curves.tex)
  % SOURCE QUOTE: Following \cite[pg.\ 1177]{V2}, let $N$ denote the rank $k-1$ kernel bundle fitting into the sequence $$0 \to N \to \mathrm{H}^0(E(-\Delta))\otimes \mathcal{O}_X \to E(-\Delta) \to 0.$$
  \textit{Source: Kemeny, Section odd-genus.}
  Define \(N\) by
  \[
    0\to N\to H^0(E(-\Delta))\otimes\cO_X\to E(-\Delta)\to0 .
  \]
  For \(t\in H^0(E)\), wedging with \(t\) gives
  \(H^0(E(-\Delta))\to H^0(L)\) and hence a map \(N\to M_L\).  Globalizing over
  \(\PP=\PP(H^0(E))\) gives
  \[
    r:N\boxtimes\cO_{\PP}(-1)\to \mathcal M=p^*M_L .
  \]
\end{definition}

\begin{definition}[Odd blow-up and elementary transform]
  \label{def:kemeny_odd_blowup_elementary_transform}
  \lean{MR4213770.OddBlowupElementaryTransform}
  \uses{def:kemeny_odd_kernel_N, lem:kemeny_decomposable}
  % SOURCE: Kemeny, Section odd-genus, elementary transformation construction (read from references/kemeny-universal-secant-bundles-syzygies-canonical-curves-source/kemeny-universal-secant-bundles-syzygies-canonical-curves.tex)
  % SOURCE QUOTE: We rectify the failure of $\mathrm{Coker}(r)$ to be locally free through a standard construction. Define $\pi: B \to X \times \PP$ as the blow-up along the codimension four locus $\mathcal{Z}$ and let $p':=p\circ \pi$, $q':=q \circ \pi$. Let $j: D \hookrightarrow B$ be the exceptional divisor.
  \textit{Source: Kemeny, Section odd-genus.}
  The map \(r\) fails to be fibrewise injective along
  \(\mathcal Z=\{(x,s)\mid s(x)=0\}\subset X\times\PP(H^0(E(-\Delta)))\).  Let
  \(\pi:B\to X\times\PP\) be the blow-up along \(\mathcal Z\), with exceptional
  divisor \(D\), and define \(p'=p\circ\pi\), \(q'=q\circ\pi\).  The elementary
  transform \(S\) is the dual of \(\operatorname{Im}((\pi^*r)^\vee)\), and
  \(\Gamma\) is the quotient in \(0\to S\to\pi^*\mathcal M\to\Gamma\to0\), where
  \(\Gamma\) is locally free of rank \(k+2\).
\end{definition}

\begin{lemma}[Odd projective-bundle model]
  \label{lem:kemeny_odd_projective_bundle_model}
  \lean{MR4213770.OddProjectiveBundleModel}
  \uses{def:kemeny_odd_blowup_elementary_transform}
  % SOURCE: Kemeny, Section odd-genus, projective-bundle model for \(B\) (read from references/kemeny-universal-secant-bundles-syzygies-canonical-curves-source/kemeny-universal-secant-bundles-syzygies-canonical-curves.tex)
  % SOURCE QUOTE: Thus $B$ is a projective bundle $\PP(\mathcal{H})$ over $\PP(\mathcal{F})$, where $\mathcal{F}:=\mathrm{Coker}(N \hookrightarrow \mathrm{H}^0(E) \otimes \mathcal{O}_X)$.
  \textit{Source: Kemeny, Section odd-genus.}
  With \(\mathcal F=\operatorname{Coker}(N\to H^0(E)\otimes\cO_X)\), the blow-up
  \(B\) is \(\PP(\mathcal H)\) over \(\PP(\mathcal F)\).  The projection
  \(\chi:B\to\PP(\mathcal F)\) satisfies
  \[
    \chi^*\cO_{\PP(\mathcal F)}(1)\simeq q'^*\cO_{\PP}(1)(-D).
  \]
\end{lemma}
\begin{proof}
  \uses{def:kemeny_odd_blowup_elementary_transform}
  Fibrewise, \(B\) is the blow-up of a projective space along the linear subspace
  cut out by \(N\).  The standard description of such a blow-up as a projective
  bundle over the quotient projective bundle gives \(\PP(\mathcal H)\), and the
  exceptional divisor accounts for the twist in \(\chi^*\cO(1)\).
\end{proof}

\begin{lemma}[Elementary transform as a relative tangent bundle]
  \label{lem:kemeny_odd_S_relative_tangent}
  \lean{MR4213770.OddSRelativeTangent}
  \uses{lem:kemeny_odd_projective_bundle_model}
  % SOURCE: Kemeny, Section odd-genus, lemma after projective-bundle model (read from references/kemeny-universal-secant-bundles-syzygies-canonical-curves-source/kemeny-universal-secant-bundles-syzygies-canonical-curves.tex)
  % SOURCE QUOTE: We have  $S \simeq T_{\chi} \otimes {q'}^*\mathcal{O}_{\PP}(-2)$, where $T_{\chi}$ is the relative tangent bundle.
  \textit{Source: Kemeny, Section odd-genus.}
  The elementary transform bundle is
  \(S\simeq T_\chi\otimes q'^*\cO_{\PP}(-2)\), where \(T_\chi\) is the relative
  tangent bundle of \(\chi:B\to\PP(\mathcal F)\).
\end{lemma}
\begin{proof}
  \uses{lem:kemeny_odd_projective_bundle_model}
  Compare the relative Euler sequence for \(\PP(\mathcal H)\) with the dual
  defining exact sequence for \(S\).  After twisting by \(q'^*\cO_{\PP}(-2)\),
  the two cokernel descriptions agree.
\end{proof}

\begin{lemma}[Determinant of \(\Gamma\) in odd genus]
  \label{lem:kemeny_odd_det_gamma}
  \lean{MR4213770.OddDetGamma}
  \uses{def:kemeny_odd_blowup_elementary_transform, lem:kemeny_odd_S_relative_tangent, def:peer_affine_pushforward_base_change}
  % SOURCE: Kemeny, Section odd-genus, Lemma det-gamma (read from references/kemeny-universal-secant-bundles-syzygies-canonical-curves-source/kemeny-universal-secant-bundles-syzygies-canonical-curves.tex)
  % SOURCE QUOTE: We have $\det \Gamma \simeq {q'}^* \mathcal{O}_{\PP}(k-1)(-D-{p'}^*\Delta)$. In particular, we have a natural isomorphism $\mathrm{H}^2(B, \wedge^{k+2} \Gamma ({p'}^*\Delta)(D))\simeq \mathrm{Sym}^{k-1}\mathrm{H}^0(X,E)^{\vee}$.
  \textit{Source: Kemeny, Lemma det-gamma.}
  One has
  \(\det\Gamma\simeq q'^*\cO_{\PP}(k-1)(-D-p'^*\Delta)\).  Consequently
  \[
    H^2\bigl(B,\wedge^{k+2}\Gamma(p'^*\Delta)(D)\bigr)
    \simeq \Sym^{k-1}H^0(X,E)^\vee .
  \]
\end{lemma}
\begin{proof}
  \uses{def:kemeny_odd_blowup_elementary_transform, lem:kemeny_odd_S_relative_tangent, def:peer_affine_pushforward_base_change}
  Taking determinants in \(0\to S\to\pi^*\mathcal M\to\Gamma\to0\) and using
  the tangent-bundle description of \(S\) gives the determinant formula.  The
  cohomology identification follows after cancelling \(D\) and \(p'^*\Delta\)
  and pushing \(q'^*\cO_{\PP}(k-1)\) down to the projective-space factor.
\end{proof}

\begin{proposition}[Most odd-genus exterior-power vanishings]
  \label{prop:kemeny_odd_most_of_vanishing}
  \lean{MR4213770.OddMostOfVanishing}
  \uses{lem:kemeny_odd_S_relative_tangent, lem:kemeny_odd_projective_bundle_model}
  % SOURCE: Kemeny, Section odd-genus, Proposition most-of-van (read from references/kemeny-universal-secant-bundles-syzygies-canonical-curves-source/kemeny-universal-secant-bundles-syzygies-canonical-curves.tex)
  % SOURCE QUOTE: We have $\mathrm{H}^{2+i}(\bigwedge^{k+2-i} \pi^* \mathcal{M} \otimes \mathrm{Sym}^i(S) ({p'}^*\Delta)(D))=0$ for $1 \leq i \leq k+2$, $i \neq k, k+1$.
  \textit{Source: Kemeny, Proposition most-of-van.}
  For \(1\le i\le k+2\) with \(i\ne k,k+1\),
  \[
    H^{2+i}\bigl(B,\wedge^{k+2-i}\pi^*\mathcal M\otimes
    \Sym^i(S)(p'^*\Delta)(D)\bigr)=0 .
  \]
\end{proposition}
\begin{proof}
  \uses{lem:kemeny_odd_S_relative_tangent, lem:kemeny_odd_projective_bundle_model}
  Resolve \(\Sym^i(S)\) by the Euler sequence for \(S\simeq T_\chi\otimes
  q'^*\cO_{\PP}(-2)\).  For \(1\le i\le k-1\), the relative projective-space
  fibres of \(\chi\), Leray, and Kunneth vanish the terms; the remaining
  \(i=k+2\) case is reduced to a Kunneth vanishing on \(X\times\PP\).
\end{proof}

\begin{lemma}[Odd Lazarsfeld--Mukai multiplication map]
  \label{lem:kemeny_LM_mult_map}
  \lean{MR4213770.LMMultMap}
  \uses{lem:kemeny_E_minus_delta_is_F, lem:kemeny_BNP_gen}
  % SOURCE: Kemeny, Section odd-genus, Lemma LM-mult-map (read from references/kemeny-universal-secant-bundles-syzygies-canonical-curves-source/kemeny-universal-secant-bundles-syzygies-canonical-curves.tex)
  % SOURCE QUOTE: The multiplication map $\mathrm{H}^0(X,E(-\Delta)) \otimes \mathrm{H}^0(X,L-\Delta) \to \mathrm{H}^0(X,E(L-2\Delta))$ is surjective.
  \textit{Source: Kemeny, Lemma LM-mult-map.}
  The multiplication map
  \[
    H^0(X,E(-\Delta))\otimes H^0(X,L-\Delta)\to H^0(X,E(L-2\Delta))
  \]
  is surjective.
\end{lemma}
\begin{proof}
  \uses{lem:kemeny_E_minus_delta_is_F, lem:kemeny_BNP_gen}
  Restrict to a smooth curve \(C\in|L-\Delta|\), split the restricted bundle as
  the two Lazarsfeld--Mukai summands associated to a pencil, and apply the
  base-point-free pencil trick, together with the standard reduction from
  restricted multiplication to the surface multiplication map.
\end{proof}

\begin{lemma}[Odd multiplication vanishings]
  \label{lem:kemeny_odd_mult_vanishings}
  \lean{MR4213770.OddMultVanishings}
  \uses{def:kemeny_odd_kernel_N, lem:kemeny_LM_mult_map}
  % SOURCE: Kemeny, Section odd-genus, Lemma mult-maps (read from references/kemeny-universal-secant-bundles-syzygies-canonical-curves-source/kemeny-universal-secant-bundles-syzygies-canonical-curves.tex)
  % SOURCE QUOTE: We have $\mathrm{H}^2(X,M_L(\Delta) \otimes N)=\mathrm{H}^2(M_L(\Delta))=0$.
  \textit{Source: Kemeny, Lemma mult-maps.}
  The vanishings
  \(H^2(X,M_L(\Delta)\otimes N)=0\) and \(H^2(X,M_L(\Delta))=0\) hold.
\end{lemma}
\begin{proof}
  \uses{def:kemeny_odd_kernel_N, lem:kemeny_LM_mult_map}
  Use the defining sequence of \(N\).  The first vanishing is reduced to
  surjectivity of the surface multiplication map, which follows from
  \(\cref{lem:kemeny_LM_mult_map}\) and the restriction diagram along
  \(\Delta\); the second follows directly from the defining sequence of \(M_L\).
\end{proof}

\begin{proposition}[Odd \(i=k+1\) vanishing]
  \label{prop:kemeny_green_odd_main}
  \lean{MR4213770.GreenOddMain}
  \uses{lem:kemeny_odd_projective_bundle_model, lem:kemeny_odd_mult_vanishings}
  % SOURCE: Kemeny, Section odd-genus, Proposition green-odd-main (read from references/kemeny-universal-secant-bundles-syzygies-canonical-curves-source/kemeny-universal-secant-bundles-syzygies-canonical-curves.tex)
  % SOURCE QUOTE: We have $\mathrm{H}^{k+3}(B,\pi^* \mathcal{M} \otimes \mathrm{Sym}^{k+1}(S) ({p'}^*\Delta)(D))=0$.
  \textit{Source: Kemeny, Proposition green-odd-main.}
  The middle vanishing
  \[
    H^{k+3}\bigl(B,\pi^*\mathcal M\otimes\Sym^{k+1}(S)(p'^*\Delta)(D)\bigr)=0
  \]
  holds.
\end{proposition}
\begin{proof}
  \uses{lem:kemeny_odd_projective_bundle_model, lem:kemeny_odd_mult_vanishings}
  Resolve \(\Sym^{k+1}(S)\) using the relative Euler sequence.  Relative duality
  over \(\chi\) turns the critical map into the multiplication map on symmetric
  powers of \(\mathcal H\), whose composite is multiplication by \(k+1\).  The
  remaining obstruction groups are exactly the Kunneth consequences of
  \(\cref{lem:kemeny_odd_mult_vanishings}\).
\end{proof}

\begin{lemma}[Odd cohomology identifications]
  \label{lem:kemeny_odd_identilem}
  \lean{MR4213770.OddIdentiLem}
  \uses{lem:kemeny_odd_projective_bundle_model}
  % SOURCE: Kemeny, Section odd-genus, Lemma identilem (read from references/kemeny-universal-secant-bundles-syzygies-canonical-curves-source/kemeny-universal-secant-bundles-syzygies-canonical-curves.tex)
  % SOURCE QUOTE: For any $j$, we have have natural isomorphisms
  \textit{Source: Kemeny, Lemma identilem.}
  For every \(j\), there are natural isomorphisms
  \[
  \begin{aligned}
    H^j(B,\pi^*\mathcal M\otimes\Sym^{k+1}(S)(p'^*\Delta)(D))
      &\simeq H^{j+1}(B,\pi^*\mathcal M\otimes S(-2k)(p'^*\Delta)(D)),\\
    H^j(B,\wedge^2\pi^*\mathcal M\otimes\Sym^k(S)(p'^*\Delta)(D))
      &\simeq H^{j+1}(B,\wedge^2\pi^*\mathcal M(-2k)(p'^*\Delta)(D)).
  \end{aligned}
  \]
\end{lemma}
\begin{proof}
  \uses{lem:kemeny_odd_projective_bundle_model}
  Apply the relative Euler resolution of \(\Sym^{k+1}(S)\) and \(\Sym^k(S)\),
  push along \(\chi\), and use relative duality.  The natural splitting of the
  symmetric-power surjection identifies the kernels and cokernels with the
  displayed cohomology groups.
\end{proof}

\begin{lemma}[Vanishing for \(\Sym^2(S)\)]
  \label{lem:kemeny_odd_sym2S_vanishing}
  \lean{MR4213770.OddSymTwoSVanishing}
  \uses{def:kemeny_odd_kernel_N, lem:kemeny_BNP_gen}
  % SOURCE: Kemeny, Section odd-genus, lemma before final proposition (read from references/kemeny-universal-secant-bundles-syzygies-canonical-curves-source/kemeny-universal-secant-bundles-syzygies-canonical-curves.tex)
  % SOURCE QUOTE: We have $\mathrm{H}^{k+4}(\mathrm{Sym}^2(S)(-2k)({p'}^*\Delta)(D))=0$.
  \textit{Source: Kemeny, Section odd-genus.}
  The cohomology group
  \(H^{k+4}(B,\Sym^2(S)(-2k)(p'^*\Delta)(D))\) vanishes.
\end{lemma}
\begin{proof}
  \uses{def:kemeny_odd_kernel_N, lem:kemeny_BNP_gen}
  Away from \(D\), \(\Sym^2(S)(p'^*\Delta)\) is identified with the pullback of
  \(\Sym^2N(\Delta)\boxtimes\cO_{\PP}(-2)\).  Kunneth reduces the claim to
  \(H^2(X,\Sym^2N(\Delta))=0\), which follows from the symmetric-square sequence
  for \(N\) and surjectivity of the determinant map, itself obtained from Petri
  generality of curves in \(|L+\Delta|\).
\end{proof}

\begin{proposition}[Odd \(\wedge^2\Gamma\)-vanishing]
  \label{prop:kemeny_odd_wedge2_gamma_vanishing}
  \lean{MR4213770.OddWedgeTwoGammaVanishing}
  \uses{lem:kemeny_odd_det_gamma, prop:kemeny_green_odd_main, lem:kemeny_odd_identilem, lem:kemeny_odd_sym2S_vanishing}
  % SOURCE: Kemeny, Section odd-genus, final proposition before odd theorem (read from references/kemeny-universal-secant-bundles-syzygies-canonical-curves-source/kemeny-universal-secant-bundles-syzygies-canonical-curves.tex)
  % SOURCE QUOTE: We have $\mathrm{H}^{k+3}(\bigwedge^2\Gamma (-2k)({p'}^*\Delta)(D))=0$. In particular, $\phi$ is surjective.
  \textit{Source: Kemeny, Section odd-genus.}
  The group
  \(H^{k+3}(B,\wedge^2\Gamma(-2k)(p'^*\Delta)(D))\) vanishes.  Consequently the
  map \(\phi\) induced by \(0\to S\to\pi^*\mathcal M\to\Gamma\to0\) is
  surjective.
\end{proposition}
\begin{proof}
  \uses{lem:kemeny_odd_det_gamma, prop:kemeny_green_odd_main, lem:kemeny_odd_identilem, lem:kemeny_odd_sym2S_vanishing}
  The four-term exterior-power sequence
  \(0\to\Sym^2(S)\to\pi^*\mathcal M\otimes S\to\wedge^2\pi^*\mathcal M
  \to\wedge^2\Gamma\to0\) reduces surjectivity to the \(\Sym^2(S)\)-vanishing
  and the \(\wedge^2\Gamma\)-vanishing.  The latter is proved by first restricting
  to \(D\), then using the determinant formula, Serre duality, and the same
  relative-Euler vanishing package as above.
\end{proof}

\begin{proposition}[Odd-genus projection-map surjectivity dual]
  \label{prop:kemeny_odd_phi_surjective}
  \lean{MR4213770.OddPhiSurjective}
  \uses{lem:kemeny_odd_det_gamma, prop:kemeny_odd_most_of_vanishing, prop:kemeny_green_odd_main, lem:kemeny_odd_identilem, prop:kemeny_odd_wedge2_gamma_vanishing}
  % SOURCE: Kemeny, Section odd-genus, propositions most-of-van, green-odd-main, and final proposition (read from references/kemeny-universal-secant-bundles-syzygies-canonical-curves-source/kemeny-universal-secant-bundles-syzygies-canonical-curves.tex)
  % SOURCE QUOTE: We have $\mathrm{H}^{k+3}(\bigwedge^2\Gamma (-2k)({p'}^*\Delta)(D))=0$. In particular, $\phi$ is surjective.
  \textit{Source: Kemeny, Section odd-genus.}
  In the odd-genus construction, the map
  \[
    \phi:H^2(\wedge^{k+2}M_L(\Delta))\to
    H^2(\wedge^{k+2}\Gamma(p'^*\Delta)(D))
  \]
  induced from \(0\to S\to\pi^*\cM\to\Gamma\to0\) is surjective.
  \end{proposition}
\begin{proof}
  \uses{lem:kemeny_odd_det_gamma, prop:kemeny_odd_most_of_vanishing, prop:kemeny_green_odd_main, lem:kemeny_odd_identilem, prop:kemeny_odd_wedge2_gamma_vanishing}
  The determinant lemma identifies the target of \(\phi\) with
  \(\Sym^{k-1}H^0(E)^\vee\).  Taking exterior powers of
  \(0\to S\to\pi^*\mathcal M\to\Gamma\to0\), twisted by \(D+p'^*\Delta\), gives
  the obstruction sequence for surjectivity.  The most-of-vanishing proposition
  kills all terms except \(i=k,k+1\); the \(i=k+1\) vanishing and the
  identifications rewrite the remaining map.  The final
  \(\wedge^2\Gamma\)-vanishing closes the last obstruction, hence \(\phi\) is
  surjective.
\end{proof}

\begin{lemma}[Odd projection comparison]
  \label{lem:kemeny_odd_projection_comparison}
  \lean{MR4213770.OddProjectionComparison}
  \uses{lem:kemeny_projection_criterion, lem:kemeny_green_pic2, prop:kemeny_odd_phi_surjective}
  % SOURCE: Kemeny, Section odd-genus, final proof comparison (read from references/kemeny-universal-secant-bundles-syzygies-canonical-curves-source/kemeny-universal-secant-bundles-syzygies-canonical-curves.tex)
  % SOURCE QUOTE: By Lemma \ref{proj-crit}, to complete the proof of Voisin's Theorem, it only remains to show that $\phi^{\vee}=pr_k \circ \psi^{\vee}$.
  \textit{Source: Kemeny, Section odd-genus.}
  With \(\psi^\vee:\Sym^{k-1}H^0(E)\simeq\K_{k,1}(X,L')\) from
  \(\cref{lem:kemeny_green_pic2}\) and \(\phi^\vee\) dual to the odd map
  \(\phi\), one has
  \[
    \phi^\vee=pr_k\circ\psi^\vee .
  \]
  Here \(L=L'-\Delta\), so the blueprint's \(\K_{k,1}(X,L')\) is the paper's
  \(\K_{k,1}(X,L+\Delta)\).
\end{lemma}
\begin{proof}
  \uses{lem:kemeny_projection_criterion, lem:kemeny_green_pic2, prop:kemeny_odd_phi_surjective}
  Blow up the codimension-two universal zero locus
  \(\widetilde{\mathcal Z}\subset X\times\PP(H^0(E))\), obtaining
  \(\widetilde B\) with exceptional divisor \(\widetilde D\).  Define
  \(\widetilde S_{L'}\) and \(\widetilde S_L\) by the two universal secant
  sequences on \(\widetilde B\).  Exterior powers identify \(\psi^\vee\) with the
  map from \(\widetilde S_{L'}\) and \(\phi^\vee\) with the corresponding map
  from \(\widetilde S_L\).  On the complement of the codimension-two locus there
  is an exact sequence
  \(0\to\widetilde S_L\to\widetilde S_{L'}\to\cO(-\widetilde p^*\Delta)\to0\),
  and the induced exterior-power diagram is exactly the projection \(pr_k\).  The
  diagram extends uniquely across codimension two, giving the equality.
\end{proof}

\begin{theorem}[Voisin's theorem in odd genus]
  \label{thm:kemeny_odd_voisin}
  \lean{MR4213770.OddVoisin}
  \uses{lem:kemeny_projection_criterion, lem:kemeny_green_pic2, prop:kemeny_odd_phi_surjective, lem:kemeny_odd_projection_comparison}
  % SOURCE: Kemeny, Section odd-genus (read from references/kemeny-universal-secant-bundles-syzygies-canonical-curves-source/kemeny-universal-secant-bundles-syzygies-canonical-curves.tex)
  % SOURCE QUOTE: In this section we prove Green's Conjecture for generic curves of odd genus $g=2k+1 \geq 9$.
  \textit{Source: Kemeny, Section odd-genus.}
  For a general complex curve \(C\) of odd genus \(g=2k+1\ge9\),
  \(\K_{k,1}(C,\omega_C)=0\).
\end{theorem}
\begin{proof}
  \uses{lem:kemeny_projection_criterion, lem:kemeny_green_pic2, prop:kemeny_odd_phi_surjective, lem:kemeny_odd_projection_comparison}
  Lemma~\ref{lem:kemeny_green_pic2} identifies the source of the projection map
  with \(\Sym^{k-1}H^0(E)\).  Surjectivity of \(\phi\) makes
  \(\phi^\vee\) injective, and the projection comparison identifies this
  injection with \(pr_k\circ\psi^\vee\).  Hence \(pr_k\) is injective.  The
  projection criterion gives \(H^1(\wedge^{k+1}M_L)=0\), and the kernel-bundle
  description identifies this with the odd-genus Koszul vanishing for a general
  canonical curve.
\end{proof}
